\documentclass[11pt]{article}
\usepackage[margin=1in]{geometry}
\usepackage{amsmath, amssymb}

\title{Exercise 2.4: General Weighting for Non-Constant Step Sizes}
\author{Sutton \& Barto, Section 2.5}
\date{}

\begin{document}
\maketitle

\section*{Problem}

If the step-size parameters, $\alpha_n$, are not constant, then the estimate $Q_n$ is a
weighted average of previously received rewards with a weighting different from that
given by (2.6). What is the weighting on each prior reward for the general case,
analogous to (2.6), in terms of the sequence of step-size parameters?

\section*{Background: Equation (2.6) for constant $\alpha$}

With a constant step size $\alpha$, the textbook derives:
\begin{equation}
    Q_{n+1} = (1-\alpha)^n Q_1 + \sum_{i=1}^{n} \alpha(1-\alpha)^{n-i}\, R_i \tag{2.6}
\end{equation}
The weight on $Q_1$ is $(1-\alpha)^n$ and the weight on $R_i$ is $\alpha(1-\alpha)^{n-i}$.
All weights sum to 1. We now generalize this to non-constant step sizes.

\section*{Solution}

\subsection*{Step 1: The update rule}

With a possibly different step size $\alpha_n$ at each step:
\[
    Q_{n+1} = Q_n + \alpha_n\bigl[R_n - Q_n\bigr] = (1 - \alpha_n)\,Q_n + \alpha_n\, R_n
\]

\subsection*{Step 2: Recursive expansion}

Expand by substituting $Q_n$, then $Q_{n-1}$, and so on:
\begin{align*}
    Q_{n+1} &= \alpha_n R_n + (1-\alpha_n)\,Q_n \\
    &= \alpha_n R_n + (1-\alpha_n)\bigl[\alpha_{n-1} R_{n-1} + (1-\alpha_{n-1})\,Q_{n-1}\bigr] \\
    &= \alpha_n R_n + (1-\alpha_n)\alpha_{n-1} R_{n-1}
       + (1-\alpha_n)(1-\alpha_{n-1})\,Q_{n-1} \\
    &\;\;\vdots
\end{align*}
Continuing until we reach $Q_1$:

\[
    \boxed{Q_{n+1} = \prod_{j=1}^{n}(1-\alpha_j)\cdot Q_1
    \;+\; \sum_{i=1}^{n}\alpha_i \prod_{j=i+1}^{n}(1-\alpha_j)\cdot R_i}
\]

where the empty product $\prod_{j=n+1}^{n}(\cdot) = 1$ (i.e., $R_n$ has weight $\alpha_n$).

\subsection*{Step 3: Identify the weights}

Define:
\begin{align*}
    w_0 &= \prod_{j=1}^{n}(1-\alpha_j)
    & &\text{(weight on initial estimate } Q_1\text{)} \\[6pt]
    w_i &= \alpha_i \prod_{j=i+1}^{n}(1-\alpha_j), \quad i = 1, \ldots, n
    & &\text{(weight on reward } R_i\text{)}
\end{align*}

So:
\[
    Q_{n+1} = w_0 \cdot Q_1 + \sum_{i=1}^{n} w_i \cdot R_i
\]

\subsection*{Step 4: Verify the weights sum to 1}

We prove $w_0 + \sum_{i=1}^{n} w_i = 1$ by induction. For $n=1$:
\[
    w_0 + w_1 = (1-\alpha_1) + \alpha_1 = 1 \;\checkmark
\]
For the inductive step, assume the weights sum to 1 after $n-1$ rewards.
The update $Q_{n+1} = (1-\alpha_n)Q_n + \alpha_n R_n$ multiplies all previous weights
by $(1-\alpha_n)$ and adds a new weight $\alpha_n$. Their sum is
$(1-\alpha_n)\cdot 1 + \alpha_n = 1$. $\square$

\subsection*{Step 5: Interpretation}

\begin{itemize}
    \item The weight on $R_i$ is $\alpha_i$ (the step size when $R_i$ was received) multiplied by
          all the ``retention factors'' $(1-\alpha_j)$ for every subsequent step $j > i$.
          Each later update retains a fraction $(1-\alpha_j)$ of the old estimate and overwrites
          a fraction $\alpha_j$ with the new reward.
    \item More recent rewards have fewer retention factors in their product, so they
          naturally receive more weight.
    \item The initial estimate $Q_1$ has no $\alpha$ factor --- it only survives through
          the retention factors. If all $\alpha_j > 0$, then $w_0 \to 0$ as $n \to \infty$
          (the initial bias vanishes).
\end{itemize}

\subsection*{Step 6: Recovering equation (2.6) as a special case}

When $\alpha_i = \alpha$ for all $i$:
\begin{align*}
    w_0 &= (1-\alpha)^n \\
    w_i &= \alpha\,(1-\alpha)^{n-i}
\end{align*}
which is exactly equation (2.6) --- the exponential recency-weighted average.

\end{document}
